% Two-page plain-English summary (finish plan C5). Numbers are macros from numbers.tex.
% Draft by Claude Code for AJ to rewrite (Decision 17).
\documentclass[11pt]{article}
\usepackage[margin=0.9in]{geometry}
\usepackage[T1]{fontenc}
\usepackage{lmodern}
\usepackage{microtype}
\usepackage{xcolor}
\usepackage{enumitem}
\setlist{nosep,leftmargin=*}
\newif\ifajfinal
\ajfinalfalse
\ifdefined\ajfinalflag\ajfinaltrue\fi
\newcommand{\ajcheck}[1]{\ifajfinal#1\else{\color{blue!60!black}#1}\textsuperscript{\color{red}[AJ]}\fi}
% Generated by gel.viz.paper_numbers from paper/tables. Do not edit.
\newcommand{\nAcctFourFortyOiBnHi}{2.6}
\newcommand{\nAcctFourFortyOiBnLo}{$-$4.2}
\newcommand{\nAcctFourFortyOiPctHi}{3.3\%}
\newcommand{\nAcctFourFortyOiPctLo}{$-$11.0\%}
\newcommand{\nAcctFourFortyWorstEps}{$-$0.68}
\newcommand{\nAcctFourFortyWorstFirm}{Oracle}
\newcommand{\nAcctFourFortyWorstFy}{2026}
\newcommand{\nAcctFourFortyWorstPct}{$-$11.0\%}
\newcommand{\nAcctMsftCapexBn}{115.9}
\newcommand{\nAcctMsftFy}{2026}
\newcommand{\nAcctMsftGpuCapexBn}{34.8}
\newcommand{\nAcctRealOiBnHi}{12.4}
\newcommand{\nAcctRealOiBnLo}{1.5}
\newcommand{\nAcctRealOiPctHi}{15.5\%}
\newcommand{\nAcctRealOiPctLo}{2.6\%}
\newcommand{\nAcctSlFourOiBnHi}{1.5}
\newcommand{\nAcctSlFourOiBnLo}{$-$4.5}
\newcommand{\nAcctSlFourOiPctHi}{1.9\%}
\newcommand{\nAcctSlFourOiPctLo}{$-$10.4\%}
\newcommand{\nAwsSpotUnchanged}{17\%}
\newcommand{\nAwsVHundredFirstYear}{2017}
\newcommand{\nAwsVHundredLastYear}{2026}
\newcommand{\nAwsVHundredPrice}{3.06}
\newcommand{\nCovCrchipHi}{76\%}
\newcommand{\nCovCrchipLo}{74\%}
\newcommand{\nCovWcrHi}{98\%}
\newcommand{\nCovWcrLo}{98\%}
\newcommand{\nCovWcuHi}{72\%}
\newcommand{\nCovWcuLo}{66\%}
\newcommand{\nEvSixEleven}{$-$0.18}
\newcommand{\nFourShape}{0.47}
\newcommand{\nFourShapeHi}{0.57}
\newcommand{\nFourShapeLo}{0.36}
\newcommand{\nFourShapeOnDemand}{$-$0.13}
\newcommand{\nFourShapeOnDemandHi}{0.13}
\newcommand{\nFourShapeOnDemandLo}{$-$0.40}
\newcommand{\nFourShapePcm}{0.29}
\newcommand{\nFourShapeSpot}{0.87}
\newcommand{\nFourShapeSpotHi}{1.41}
\newcommand{\nFourShapeSpotLo}{0.34}
\newcommand{\nFourShapeWcrHi}{1.03}
\newcommand{\nFourShapeWcrLo}{$-$0.44}
\newcommand{\nFourShapeWcrP}{0.29}
\newcommand{\nFourShapeWcuP}{$<$0.0001}
\newcommand{\nGoogleShare}{62\%}
\newcommand{\nHHundredDispersion}{3.52}
\newcommand{\nHHundredHypNeo}{3.10}
\newcommand{\nHTwo}{$-$0.09}
\newcommand{\nHTwoHi}{0.23}
\newcommand{\nHTwoLo}{$-$0.40}
\newcommand{\nHyperscalerShare}{94\%}
\newcommand{\nListMoversRows}{5{,}226}
\newcommand{\nListTransitions}{9{,}282}
\newcommand{\nListUnchanged}{96\%}
\newcommand{\nMarkupShare}{47\%}
\newcommand{\nMcReps}{200}
\newcommand{\nNeocloudRows}{1{,}538}
\newcommand{\nOneARateEarly}{0.27}
\newcommand{\nOneARateEarlyHi}{0.77}
\newcommand{\nOneARateEarlyLo}{$-$0.23}
\newcommand{\nOneARateEarlyNoHTwoHundred}{$-$0.11}
\newcommand{\nOneARateLate}{$-$0.03}
\newcommand{\nOneARateLateHi}{0.18}
\newcommand{\nOneARateLateLo}{$-$0.24}
\newcommand{\nOneAShape}{0.31}
\newcommand{\nOneAShapeHi}{1.11}
\newcommand{\nOneAShapeLo}{$-$0.52}
\newcommand{\nOneBRateEarly}{0.28}
\newcommand{\nOneBRateEarlyHi}{0.84}
\newcommand{\nOneBRateEarlyLo}{$-$0.28}
\newcommand{\nOneBRateLate}{0.00}
\newcommand{\nOneBRateLateHi}{0.47}
\newcommand{\nOneBRateLateLo}{$-$0.50}
\newcommand{\nOneBShape}{0.28}
\newcommand{\nOneBShapeHi}{0.97}
\newcommand{\nOneBShapeLo}{$-$0.42}
\newcommand{\nPOneDHHundred}{$-$0.02}
\newcommand{\nPThreeAbundant}{11.0}
\newcommand{\nPThreeMedian}{4.1}
\newcommand{\nPThreeNoScarcity}{10.7}
\newcommand{\nPThreePNinety}{7.0}
\newcommand{\nPThreePTen}{2.2}
\newcommand{\nPTwoLopexHi}{11.7}
\newcommand{\nPTwoLopexLo}{9.8}
\newcommand{\nPanelRows}{24{,}210}
\newcommand{\nPathAHundredFive}{$-$0.45}
\newcommand{\nPathAHundredFivePct}{36\%}
\newcommand{\nPathHHundredFour}{$-$0.04}
\newcommand{\nPathHHundredFourPct}{4\%}
\newcommand{\nPathHHundredOne}{0.24}
\newcommand{\nPathHHundredOnePct}{$-$27\%}
\newcommand{\nPathHHundredThree}{$-$0.50}
\newcommand{\nPathHHundredThreePct}{40\%}
\newcommand{\nPathVHundredFour}{$-$0.13}
\newcommand{\nPathVHundredFourPct}{12\%}
\newcommand{\nPathVHundredNine}{$-$0.38}
\newcommand{\nPathVHundredNinePct}{32\%}
\newcommand{\nPermAssignments}{1{,}260}
\newcommand{\nPermP}{0.0008}
\newcommand{\nPermSharePositive}{55\%}
\newcommand{\nRawRows}{61{,}873}
\newcommand{\nSampleRows}{20{,}185}
\newcommand{\nShapeChipTrendZero}{$-$0.39}
\newcommand{\nShapeEqualProv}{0.19}
\newcommand{\nShapeListMovers}{$-$0.52}
\newcommand{\nShapeListMoversSe}{0.24}
\newcommand{\nShapeNoGoogle}{0.44}
\newcommand{\nTflopsCoef}{0.88}
\newcommand{\nTstarAll}{never}
\newcommand{\nTwoRateEarly}{0.11}
\newcommand{\nTwoRateEarlyHi}{0.35}
\newcommand{\nTwoRateEarlyLo}{$-$0.12}
\newcommand{\nTwoRateLate}{$-$0.15}
\newcommand{\nTwoRateLateHi}{0.11}
\newcommand{\nTwoRateLateLo}{$-$0.41}
\newcommand{\nTwoShape}{0.26}
\newcommand{\nTwoShapeHi}{0.47}
\newcommand{\nTwoShapeLo}{0.05}
\newcommand{\nTwoShapeWcrP}{0.30}
\newcommand{\nValFourFortyLostOne}{27\%}
\newcommand{\nValFourFortyLostThree}{72\%}
\newcommand{\nValFourFortyRentZero}{16.02}
\newcommand{\nValFourFortySle}{4.9}
\newcommand{\nValFourFortyVzeroCapex}{8.9}
\newcommand{\nValFourTwentyLostOne}{15\%}
\newcommand{\nValFourTwentyLostThree}{53\%}
\newcommand{\nValFourTwentyRentZero}{6.96}
\newcommand{\nValFourTwentySle}{7.1}
\newcommand{\nValFourTwentyVzeroCapex}{5.9}
\newcommand{\nValOpex}{0.30}
\newcommand{\nValRealLostOne}{11\%}
\newcommand{\nValRealLostThree}{32\%}
\newcommand{\nValRealRentZero}{5.85}
\newcommand{\nValRealSle}{9.4}
\newcommand{\nValRealSleFive}{5.1}
\newcommand{\nValRealSleSix}{6.1}
\newcommand{\nValRealVzeroCapex}{4.5}
\newcommand{\nValRealVzeroHi}{6.3}
\newcommand{\nValRealVzeroLo}{2.9}
\newcommand{\nValRentNow}{4.50}
\newcommand{\nValTwoLostOne}{6\%}
\newcommand{\nValTwoLostThree}{28\%}
\newcommand{\nValTwoRentZero}{4.83}
\newcommand{\nValTwoSle}{10.1}
\newcommand{\nValTwoVzeroCapex}{4.6}

\pagestyle{empty}
\begin{document}
\begin{center}
{\Large\bfseries How fast do AI chips lose value?}\\[2pt]
{\large What nine years of GPU rental prices say about tech-company depreciation}\\[4pt]
Ajai Upadhyaya $\cdot$ two-page summary of the working paper ``Rental-Implied Depreciation of AI Accelerators''
\end{center}

\paragraph{The question.} The biggest AI spenders (Microsoft, Alphabet, Amazon, Meta, Oracle,
CoreWeave, Nebius) spread the cost of their servers over five to six years. If AI chips actually
lose their value faster, today's profits are overstated; if slower, understated. Published
estimates of how fast the chips lose value disagree wildly, from roughly flat to about 40\% a year.

\paragraph{What I did.} I built an open dataset of \nPanelRows{} monthly rental prices, centered on
seven NVIDIA data-center chips from the 2017 V100 to the 2024 B200, across cloud
providers and contract types, mostly from public price files and archived web pages. I then asked
how a chip's rent changes as it ages, turned rents into asset values, and compared the result
with the companies' accounting lives.

\paragraph{The hard part.} A chip's age goes up exactly as fast as the calendar, so ``older chips
rent for less'' and ``rents fell this year'' are hard to tell apart. I report what actually
happened to each chip and four different ways of separating age from time, side by side, and I
had an independent reviewer try to break the analysis before writing it up.

\paragraph{What I found.}
\begin{itemize}
\item \textbf{Rents fall slowly.} Holding the seller and contract fixed, the 2017 V100's rent
fell by about \nPathVHundredNinePct{} over nine years.
\item \textbf{Most posted prices hardly ever move.} \nListUnchanged{} of month-to-month on-demand
and reserved price observations show no change. Amazon has charged \$\nAwsVHundredPrice{} an hour
for a V100 since \nAwsVHundredFirstYear{}.
\item \textbf{Which market you look at changes the answer.} Spot prices fall late in a chip's
life; list prices that do move fall early. Published estimates disagree partly because they read
different markets.
\item \textbf{With honest statistics, the data cannot pin down a precise rate.} With only seven
chips, standard methods overstate precision: in a simulation they caught the true value
only \nCovWcuLo--\nCovWcuHi{} of the time instead of 95\%. Corrected, no method finds a
statistically clear decline.
\item \textbf{Rents never get close to running costs.} An H100 rents for about
\$\nValRentNow{} an hour against about \$\nValOpex{} of electricity and maintenance. So a chip is
not retired because it stops covering its costs; it is retired when a newer chip would earn more
in the same scarce power slot. A retirement model puts that at about \nPThreeMedian{} years, and
the companies' five- to six-year book lives sit near it.
\end{itemize}

\paragraph{Why it matters.} \ajcheck{The depreciation debate is less about facts than about three
unstated choices: which price to read, how to separate age from time, and what ends a chip's
useful life. Making those choices explicit, with open data and code, is the paper's contribution.}

\paragraph{Skills used.} Building a panel dataset from APIs, price files and web archives (Python,
Polars, DuckDB); fixed-effects regression and few-cluster inference (wild bootstrap, Monte Carlo
coverage checks); asset valuation; SEC XBRL financial data; a fully reproducible pipeline
(\texttt{make all} rebuilds every number from raw data).

{\footnotesize The code, analysis and drafts were produced with Claude (Anthropic) in Claude Code
under the author's direction; the author set the questions and approved the results.}
\end{document}
